\documentclass[11pt]{article}
\usepackage{amsmath,amssymb,amsthm}
\usepackage{booktabs}
\usepackage[margin=1in]{geometry}

\newtheorem{theorem}{Theorem}
\newtheorem{lemma}[theorem]{Lemma}

\title{Problem 681: Maximal Area}
\author{Project Euler Solutions}
\date{}

\begin{document}
\maketitle

\section{Problem Statement}

For positive integers $a \le b \le c \le d$, let $M(a,b,c,d)$ denote the maximum area of a quadrilateral with consecutive side lengths $a, b, c, d$. Define
\[
SP(n) = \sum_{\substack{a \le b \le c \le d \\ M(a,b,c,d) \in \{1,\ldots,n\}}} (a + b + c + d).
\]
Given: $SP(10) = 186$, $SP(100) = 23238$. Find $SP(10^{16})$.

\section{Mathematical Foundation}

\begin{theorem}[Brahmagupta]
Among all quadrilaterals with prescribed side lengths $a, b, c, d > 0$ satisfying $d < a+b+c$, the maximum area is attained uniquely by the cyclic quadrilateral and equals
\[
M(a,b,c,d) = \sqrt{(s-a)(s-b)(s-c)(s-d)},
\]
where $s = (a+b+c+d)/2$ is the semi-perimeter.
\end{theorem}

\begin{proof}
For a quadrilateral with sides $a,b,c,d$ and opposite angle pairs $(\alpha,\gamma)$, $(\beta,\delta)$, the area satisfies
\[
A = \sqrt{(s-a)(s-b)(s-c)(s-d) - abcd\cos^2\!\Bigl(\frac{\alpha+\gamma}{2}\Bigr)}.
\]
Since $\cos^2((\alpha+\gamma)/2) \ge 0$ with equality iff $\alpha+\gamma=\pi$ (the cyclic condition), the maximum is $\sqrt{(s-a)(s-b)(s-c)(s-d)}$.
\end{proof}

\begin{lemma}[Change of Variables]
Define $w = s-a$, $x = s-b$, $y = s-c$, $z = s-d$. Then:
\begin{enumerate}
  \item $a+b+c+d = w+x+y+z$ and $s = (w+x+y+z)/2$.
  \item $a \le b \le c \le d \iff w \ge x \ge y \ge z$.
  \item The polygon inequality $d < a+b+c$ holds iff $z > 0$.
  \item $M(a,b,c,d) \in \mathbb{Z}^+$ requires $w+x+y+z$ even and $wxyz$ a perfect square.
\end{enumerate}
\end{lemma}

\begin{proof}
We have $w+x+y+z = 4s - 2s = 2s = a+b+c+d$. Part~(2): $a \le b$ iff $s-w \le s-x$ iff $w \ge x$. Part~(3): $d < a+b+c$ iff $s-z < 3s-(w+x+y)$, simplifying (via $w+x+y+z=2s$) to $z > 0$. Part~(4): $M \in \mathbb{Z}^+$ requires $s \in \mathbb{Z}$ (even perimeter) and $wxyz$ a perfect square.
\end{proof}

\begin{lemma}[Enumeration Bijection]
There is a bijection between valid quadrilateral tuples $(a,b,c,d)$ with integer $M = A$ and factorizations $A^2 = wxyz$ with $w \ge x \ge y \ge z \ge 1$ and $w+x+y+z$ even.
\end{lemma}

\begin{proof}
Given such a factorization, set $s=(w+x+y+z)/2$ and recover $(a,b,c,d)=(s-w,s-x,s-y,s-z)$. By the previous lemma, this yields a valid ordered quadrilateral with $M=A$. The map is invertible.
\end{proof}

\section{Algorithm}

Enumerate all 4-fold factorizations of $A^2$ for each $A = 1,\ldots,n$:
\begin{enumerate}
  \item For $z = 1$ to $A^{1/2}$: if $z \mid A^2$, set $R_1 = A^2/z$.
  \item For $y = z$ to $R_1^{1/3}$: if $y \mid R_1$, set $R_2 = R_1/y$.
  \item For $x = y$ to $R_2^{1/2}$: if $x \mid R_2$, set $w = R_2/x$.
  \item If $w+x+y+z$ is even, accumulate $w+x+y+z$ into $SP(n)$.
\end{enumerate}

\section{Complexity Analysis}

\begin{itemize}
  \item \textbf{Time:} $O\bigl(\sum_{A=1}^{n} d_4(A^2)\bigr) = O(n\log^3 n)$ on average.
  \item \textbf{Space:} $O(n)$ with sieve, or $O(\sqrt{n})$ with trial division.
\end{itemize}

\section{Answer}

\[
\boxed{2611227421428}
\]

\end{document}
